\documentclass[10pt,a4paper]{amsart}
\usepackage[margin=19mm]{geometry}
\usepackage{amsmath,amssymb,mathtools}
\usepackage[scr=rsfs,cal=euler]{mathalfa}
\usepackage{xfrac}
\usepackage[unicode,hidelinks,bookmarks=false]{hyperref}
\newcommand{\ie}{i.e.\ }
\newcommand{\cf}{cf.\ }
\newcommand{\resp}{respectively\ }
\newcommand{\Subsection}{\subsection}
\providecommand{\nobreakdash}{\nobreak\hbox{-}}
\setlength{\emergencystretch}{2em}
\multlinegap0pt
\usepackage{amssymb}
\usepackage{mathtools}
\usepackage{xcolor}
\newtheorem{theorem}{Theorem}
\title{One-Cut Risk Profiles under Quadratic Loss:\\ Discrete Convexity, Continuous Limits, and Higher Dimensions}
\author{Mihaela-Adriana Nistor, Ionel Popescu}
\date{Source: arXiv:2609.05357v1 (CC BY 4.0)}
\begin{document}
\maketitle

\noindent\emph{Source and licence.} One-Cut Risk Profiles under Quadratic Loss:\\ Discrete Convexity, Continuous Limits, and Higher Dimensions. Version arXiv:2609.05357v1 (CC BY 4.0), distributed by arXiv under the Creative Commons Attribution 4.0 International licence, http://creativecommons.org/licenses/by/4.0/, as recorded in the arXiv licence metadata for this exact version and shown on the abstract page https://arxiv.org/abs/2609.05357v1. Full licence notice: \texttt{licenses/arxiv-2609.05357-CC-BY-4.0.txt}.\par\smallskip

\noindent\textit{Editorial scope.} Counted unit: the article's theorem thm:fiber-reduction (source0.tex lines 872--903). The article's complete original proof of it is delivered with it (source0.tex lines 905--987, one \texttt{proof} environment of 83 lines). No mathematics is written, repaired or re-derived: the statement and the proof are the article's own lines, in the article's own notation, and no retained line is substituted or re-flowed.  No further source-local context is needed: every cross-reference of every retained line resolves inside the two delivered spans.  Nothing was omitted from any retained span, so the package carries no omission record.  No retained line contains a citation, so the wrapper carries no bibliography and no bibliography entry.  No retained line contains a figure environment, an image-inclusion command or drawing code, so no image is delivered and the package declares no asset dependency.  Editorial formatting only, in addition: an AMS \texttt{amsart} preamble that restates the article's own declarations and macros used by the retained lines, namely the environments \texttt{theorem} on separate counters, together with the standard packages those lines need.  The retained environments print in this document as Theorem 1 in order of appearance, whereas the article prints the same items with the numbering of its own full document.  The complete native source of the article as distributed in the arXiv e-print for this exact version is bundled unchanged as \texttt{sources/209427/source0.tex}.
\begin{theorem}[Fixed-direction fiber reduction]
\label{thm:fiber-reduction}
Let $z(Y):=z_i$ on $\{Y=t_i\}$ and
\[
  \sigma^2_{\mathrm{fib}}
  :=\mathbb E\|X-\mu-z(Y)\|_M^2.
\]
Then, for $1\leq i\leq m-1$,
\begin{equation}\label{eq:fiber-decomposition}
  V_i=\sigma^2_{\mathrm{fib}}
      +\sum_{j\leq i}p_j\|z_j-a_i\|_M^2
      +\sum_{j>i}p_j\|z_j-b_i\|_M^2.
\end{equation}
At the endpoints,
\[
 V_0=V_m=T=\sigma^2_{\mathrm{fib}}+\sum_{j=1}^m p_j\|z_j\|_M^2.
\]
Thus the dispersion inside the fibers contributes the same constant to every
cut.  For $1\leq i\leq m-1$, put
\[
  S_i:=\mathbb E[(X-\mu)\mathbf 1_{\{Y\leq t_i\}}]
      =P_i a_i=-q_i b_i,
\]
then
\begin{equation}\label{eq:vector-gain}
  V_i=T-B_i,
  \qquad
  B_i=P_iq_i\|a_i-b_i\|_M^2
     =\frac{\|S_i\|_M^2}{P_iq_i},
\end{equation}
with $B_0=B_m=0$.
\end{theorem}

\begin{proof}
Set $\varepsilon:=X-\mu-z(Y)$.  We have
$\mathbb E[\varepsilon\mid Y]=0$.  If $c$ is any vector-valued function of
$Y$, then
\[
 \mathbb E\langle\varepsilon,c(Y)\rangle_M
 =\mathbb E\left[
   \left\langle\mathbb E[\varepsilon\mid Y],c(Y)\right\rangle_M
  \right]=0.
\]
Thus the part of $X-\mu$ orthogonal to the fibers has zero cross term with
every quantity which depends only on the fiber.

For the cut after $t_i$, the conditional mean on the left is
\[
 \mathbb E[X\mid Y\leq t_i]
 =\mu+\mathbb E[z(Y)\mid Y\leq t_i]
 =\mu+a_i,
\]
and similarly the conditional mean on the right is $\mu+b_i$.  On the left
side we can therefore write
\[
 X-(\mu+a_i)=\varepsilon+z(Y)-a_i.
\]
Expanding the squared $M$--norm and conditioning on $Y$ gives
\begin{align*}
 &\mathbb E\left[
   \|X-(\mu+a_i)\|_M^2\mathbf 1_{\{Y\leq t_i\}}
  \right]\\
 &\qquad=
 \mathbb E\left[\|\varepsilon\|_M^2
                 \mathbf 1_{\{Y\leq t_i\}}\right]
 +\sum_{j\leq i}p_j\|z_j-a_i\|_M^2,
\end{align*}
because the cross term vanishes.  The same calculation on the right yields
\[
 \mathbb E\left[
   \|X-(\mu+b_i)\|_M^2\mathbf 1_{\{Y>t_i\}}
  \right]
 =\mathbb E\left[\|\varepsilon\|_M^2
                 \mathbf 1_{\{Y>t_i\}}\right]
 +\sum_{j>i}p_j\|z_j-b_i\|_M^2.
\]
Adding the two identities proves \eqref{eq:fiber-decomposition}.  Taking one
regime only, with representative $\mu$, gives
\[
 T=\sigma^2_{\mathrm{fib}}+sum_{j=1}^m p_j\|z_j\|_M^2.
\]

It remains to identify the gain.  The usual within/between decomposition for
the fiber barycenters gives
\begin{align*}
 \sum_{j=1}^m p_j\|z_j\|_M^2
 &=\sum_{j\leq i}p_j\|z_j-a_i\|_M^2
   +\sum_{j>i}p_j\|z_j-b_i\|_M^2\\
 &\quad+P_i\|a_i\|_M^2+q_i\|b_i\|_M^2.
\end{align*}
Consequently,
\[
 V_i=T-\left(P_i\|a_i\|_M^2+q_i\|b_i\|_M^2\right).
\]
Since the barycenters are centered,
\[
 P_i a_i+q_i b_i=\sum_{j=1}^m p_jz_j
 =\mathbb E(X-\mu)=0.
\]
Using this identity, we obtain
\[
 P_i\|a_i\|_M^2+q_i\|b_i\|_M^2
 =P_iq_i\|a_i-b_i\|_M^2.
\]
Finally, $S_i=P_i a_i=-q_i b_i$, and hence
\[
 a_i-b_i=\frac{S_i}{P_i}+\frac{S_i}{q_i}
          =\frac{S_i}{P_iq_i}.
\]
Therefore
\[
 P_iq_i\|a_i-b_i\|_M^2
 =\frac{\|S_i\|_M^2}{P_iq_i},
\]
which proves all the identities in \eqref{eq:vector-gain}.
\end{proof}

\end{document}
